\documentclass[12pt,a4paper]{article}
\usepackage[utf8]{inputenc}
\usepackage[T1]{fontenc}
\usepackage[ngerman]{babel}
\usepackage{microtype}

\usepackage{amsmath, amssymb, amsthm}
\usepackage{physics}
\usepackage{siunitx}
\usepackage{tikz}
\usepackage{tikz-3dplot}
\usetikzlibrary{arrows.meta, decorations.pathmorphing, decorations.markings,
                patterns, calc, positioning, shapes.geometric, angles, quotes,
                decorations.pathreplacing, backgrounds, hobby, shadings, fadings}
\usepackage{pgfplots}
\pgfplotsset{compat=1.18}
\usepackage{booktabs}
\usepackage{array}
\usepackage{multirow}
\usepackage{geometry}
\usepackage{graphicx}
\usepackage{xcolor}
\usepackage{hyperref}
\usepackage{cleveref}
\usepackage{float}
\usepackage{caption}
\usepackage{subcaption}
\usepackage{enumitem}
\usepackage{fancyhdr}
\usepackage{titlesec}
\usepackage{mathtools}
\usepackage{bm}
\usepackage{rotating}

\geometry{
  a4paper,
  left=2.5cm, right=2.5cm,
  top=3cm, bottom=3cm
}

% ---- Farben ----
\definecolor{raumblau}{RGB}{10,40,100}
\definecolor{triebwerkorange}{RGB}{220,100,20}
\definecolor{hellgrau}{RGB}{240,240,245}
\definecolor{dunkelgrau}{RGB}{80,80,90}
\definecolor{flammenrot}{RGB}{200,50,20}
\definecolor{exhaust}{RGB}{180,200,230}
\definecolor{hyprot}{RGB}{180,30,20}
\definecolor{hypblau}{RGB}{15,50,120}
\definecolor{schockgelb}{RGB}{220,160,10}
\definecolor{plasmaviolett}{RGB}{120,30,160}
\definecolor{ablgruen}{RGB}{30,120,60}
\definecolor{hitzeorange}{RGB}{210,90,20}
\definecolor{wobbelviolett}{RGB}{100,20,140}
\definecolor{wobbelgruen}{RGB}{20,110,60}

\hypersetup{
  colorlinks=true,
  linkcolor=raumblau,
  citecolor=raumblau,
  urlcolor=raumblau,
  pdftitle={Raketenantriebe, Bahnmechanik und Hyperschallaerodynamik},
  pdfauthor={Stephan Epp}
}

\newtheorem{satz}{Satz}[section]
\newtheorem{definition}{Definition}[section]
\newtheorem{herleitung}{Herleitung}[section]
\newtheorem{bemerkung}{Bemerkung}[section]

% ---- Makros ----
\newcommand{\vek}[1]{\boldsymbol{#1}}
\newcommand{\tensor}[1]{\underline{\underline{#1}}}
\newcommand{\Isp}{I_{\mathrm{sp}}}
\newcommand{\ve}{v_e}
\newcommand{\mdot}{\dot{m}}
\newcommand{\Dv}{\Delta v}
\newcommand{\Ma}{M\!a}
\newcommand{\Rey}{Re}

%-------------------------------------------------------
\title{Raketenantriebe, Bahnmechanik und Hyperschallaerodynamik\\[0.4em]
  \large\textit{Formale Grundlagen, Antriebskonzepte, Satellitenbahnmechanik\\
und atmosph\"{a}risches Hyperschallregime}}
\author{Stephan Epp}
\date{\today}

\begin{document}

\maketitle

\tableofcontents
\newpage

%=======================================================
\section{Einleitung und physikalische Grundlagen}
%=======================================================

Die Raumfahrt steht und fällt mit einer einzigen fundamentalen Frage: \emph{Wie bewegt
man sich in einem Medium fort, das keinen Widerstand bietet?} Der klassische
Kolbenmotor überträgt Kraft auf Boden oder Luft; im Vakuum des Weltraums ist beides
nicht vorhanden. Die Lösung liefert das dritte Newtonsche Axiom in seiner reinsten
Form: \textit{actio} = \textit{reactio}. Jede Masse, die mit hoher Geschwindigkeit
nach hinten ausgestoßen wird, erzeugt einen Impuls nach vorne --- dies ist das
Raketenprinzip.

Dieser Aufsatz vereint drei Gebiete: die \textbf{Antriebslehre} mit Tsiolkowski-Gleichung,
Lavaldüse und verschiedenen Triebwerkskonzepten; die \textbf{Bahnmechanik} mit Hohmann-Transfer,
Satellitenorbits und einem vollständig hergeleiteten mechanischen \emph{Wobble-Effekt},
der eine Rakete gegen aktive Abwehrsysteme resistent macht; sowie die
\textbf{Hyperschallaerodynamik} für den atmosphärischen Wiedereintritt.

\subsection{Das Raketenprinzip --- Impulserhaltung}

Betrachten wir eine Rakete der Masse $m(t)$ im massefreien Raum. Im Zeitintervall
$\dd t$ wird eine Treibstoffmasse $|\dd m|$ mit der Ausströmgeschwindigkeit $\ve$
relativ zur Rakete ausgestoßen.

\begin{herleitung}[Raketengrundgleichung (Tsiolkowski)]
\textbf{Impuls zum Zeitpunkt $t$:}
\[
  p(t) = m(t)\,v(t)
\]
\textbf{Impuls zum Zeitpunkt $t + \dd t$:}
\[
  p(t+\dd t) = \underbrace{(m+\dd m)(v+\dd v)}_{\text{Rakete}} +
               \underbrace{(-\dd m)(v - \ve)}_{\text{ausgestoßenes Gas}}
\]
(da $\dd m < 0$ ist $-\dd m > 0$ die ausgestoßene Masse)

\textbf{Impulserhaltung} ($\Delta p_{\text{ext}} = 0$ im kräftefreien Raum):
\begin{align}
  p(t+\dd t) &= p(t) \notag\\
  (m+\dd m)(v+\dd v) + (-\dd m)(v-\ve) &= mv \notag\\
  mv + m\,\dd v + v\,\dd m + \dd m\,\dd v - v\,\dd m + \ve\,\dd m &= mv \notag
\end{align}
Vernachlässigung des Produkts $\dd m\,\dd v$ (Terme zweiter Ordnung):
\[
  \boxed{m\,\dd v = -\ve\,\dd m}
\]
Dies ist die \textbf{Raketengrundgleichung}.
\end{herleitung}

%=======================================================
\section{Tsiolkowski-Gleichung und Stufungsoptimierung}
%=======================================================

\subsection{Die Tsiolkowski-Gleichung}

Die Integration der Raketengrundgleichung liefert die berühmte
Tsiolkowski-Gleichung:

\begin{align}
  \int_{v_0}^{v_1} \dd v &= -\ve \int_{m_0}^{m_1} \frac{\dd m}{m} \notag\\
  v_1 - v_0 &= \ve \ln\!\left(\frac{m_0}{m_1}\right)
\end{align}

\begin{satz}[Tsiolkowski-Gleichung]
Das erreichbare Geschwindigkeitsinkrement $\Dv$ einer Rakete im kräftefreien Raum
beträgt:
\[
  \boxed{\Dv = \ve \ln\!\left(\frac{m_0}{m_1}\right) = \ve \ln\!\left(\frac{m_0}{m_0 - m_p}\right)}
\]
wobei $m_0$ die Startmasse, $m_1$ die Endmasse und $m_p = m_0 - m_1$ die
verbrauchte Treibstoffmasse bezeichnet.
\end{satz}

\begin{figure}[H]
\centering
\begin{tikzpicture}[scale=1.0]
\begin{axis}[
  width=13cm, height=8cm,
  xlabel={Massenverhältnis $R = m_0/m_1$},
  ylabel={$\Dv / \ve$},
  xmin=1, xmax=20,
  ymin=0, ymax=3.5,
  grid=both,
  grid style={line width=0.3pt, draw=gray!30},
  major grid style={line width=0.6pt, draw=gray!50},
  tick label style={font=\small},
  label style={font=\small\bfseries},
  title style={font=\normalsize\bfseries\color{raumblau}},
  title={Tsiolkowski: $\Dv/\ve = \ln(R)$ in Abhängigkeit des Massenverhältnisses},
  legend pos=north west,
  legend style={font=\small},
]
\addplot[thick, color=raumblau, domain=1:20, samples=200] {ln(x)};
\addlegendentry{$\Dv/\ve = \ln(R)$}
\addplot[only marks, mark=*, mark size=3pt, color=triebwerkorange]
  coordinates {(7.39, 2.0) (2.72, 1.0) (20.09, 3.0)};
\addlegendentry{typische Missionen}
\draw[dashed, gray] (axis cs:7.39,0) -- (axis cs:7.39,2.0)
  node[above, font=\scriptsize] {$R\approx 7.4$};
\draw[dashed, gray] (axis cs:1,2.0) -- (axis cs:7.39,2.0);
\node[font=\scriptsize, color=dunkelgrau] at (axis cs:10,2.2)
  {LEO-Mission ($\Dv=9.4$ km/s)};
\draw[dashed, gray!60] (axis cs:2.72,0) -- (axis cs:2.72,1.0);
\draw[dashed, gray!60] (axis cs:1,1.0) -- (axis cs:2.72,1.0);
\node[font=\scriptsize, color=dunkelgrau] at (axis cs:5.5,0.7)
  {$R=e \approx 2.72$};
\end{axis}
\end{tikzpicture}
\caption{Tsiolkowski-Diagramm: normiertes Geschwindigkeitsinkrement als Funktion
des Massenverhältnisses. Ein $\Dv$ von $\SI{9.4}{km/s}$ erfordert bei
$\ve = \SI{4.4}{km/s}$ ein Massenverhältnis von etwa 8.}
\label{fig:tsiol}
\end{figure}

\subsection{Stufung --- Mehrere Raketenstufen}

Da das Massenverhältnis durch Strukturmasse begrenzt ist, verwendet man \emph{mehrstufige}
Raketen. Der Gesamtgewinn durch $n$ Stufen:
\[
  \Dv_{\mathrm{ges}} = \sum_{k=1}^{n} v_{e,k}\,\ln\!\left(\frac{m_{0,k}}{m_{1,k}}\right)
\]

\begin{figure}[H]
\centering
\begin{tikzpicture}[scale=0.85]
  \begin{axis}[
    width=13cm, height=9cm,
    xlabel={Strukturmassenanteil $\varepsilon = m_s/(m_s+m_p)$},
    ylabel={$\Dv_{\mathrm{ges}}/v_e$},
    xmin=0.02, xmax=0.25,
    ymin=0, ymax=8,
    grid=both,
    grid style={draw=gray!25},
    title={Mehrstufenrakete: $\Dv$-Gewinn durch Stufenanzahl},
    title style={font=\normalsize\bfseries\color{raumblau}},
    legend pos=north east,
    legend style={font=\small},
    tick label style={font=\small},
  ]
  \addplot[thick, raumblau, domain=0.02:0.25, samples=100]
    {ln((1-x)/x)};
  \addlegendentry{1 Stufe}
  \addplot[thick, triebwerkorange, domain=0.02:0.25, samples=100]
    {2*ln((1-x)/x)};
  \addlegendentry{2 Stufen}
  \addplot[thick, red!70!black, domain=0.02:0.25, samples=100]
    {3*ln((1-x)/x)};
  \addlegendentry{3 Stufen}
  \addplot[dashed, gray, domain=0:8] coordinates {(0.07,0)(0.07,8)};
  \node[font=\scriptsize, gray, rotate=90] at (axis cs:0.075,5.5)
    {Saturn~V $(\varepsilon\approx0.07)$};
  \end{axis}
\end{tikzpicture}
\caption{Gewinn durch Stufenanzahl: Bei einem typischen Strukturanteil $\varepsilon = 0.07$
und $v_e = \SI{4.4}{km/s}$ ermöglicht eine dreistufige Rakete
$\Dv \approx \SI{33}{km/s}$.}
\label{fig:stufen}
\end{figure}

%=======================================================
\section{Thermodynamische Grundlagen des Raketentriebwerks}
%=======================================================

\subsection{De-Laval-Düse und isentrope Strömung}

Das Herzstück jedes chemischen Raketentriebwerks ist die \emph{konvergent-divergente
Lavaldüse}. Sie wandelt die thermische und chemische Energie des Treibgases in
gerichtete kinetische Energie um.

\begin{figure}[H]
\centering
\begin{tikzpicture}[scale=1.2]
  \def\Rc{1.2}
  \def\Rt{0.35}
  \def\Re{0.9}
  \draw[thick, raumblau, line width=1.5pt]
    (-3.5, \Rc) -- (-1.8, \Rc)
    .. controls (-0.8, \Rc) and (-0.3, \Rt+0.15) .. (0, \Rt)
    .. controls (0.5, \Rt-0.05) and (1.2, \Re-0.3) .. (3.0, \Re);
  \draw[thick, raumblau, line width=1.5pt]
    (-3.5,-\Rc) -- (-1.8,-\Rc)
    .. controls (-0.8,-\Rc) and (-0.3,-\Rt-0.15) .. (0,-\Rt)
    .. controls (0.5,-\Rt+0.05) and (1.2,-\Re+0.3) .. (3.0,-\Re);
  \draw[thick, raumblau, line width=1.5pt] (-3.5,-\Rc) -- (-3.5,\Rc);
  \draw[thick, raumblau, line width=1.5pt] (3.0,-\Re) -- (3.0,\Re);
  \foreach \yf in {-0.9,-0.6,-0.3,0,0.3,0.6,0.9} {
    \pgfmathsetmacro{\len}{0.4}
    \fill[triebwerkorange, opacity=0.25]
      (-3.0,\yf-0.15) .. controls (-2.0+\len,\yf) .. (-1.5,\yf+0.15)
      .. controls (-1.5,\yf) .. (-3.0,\yf+0.05) -- cycle;
  }
  \foreach \y in {-0.8,-0.4,0,0.4,0.8} {
    \draw[-{Stealth}, exhaust!80!blue, line width=0.8pt, opacity=0.8]
      (-1.5, \y*0.9) .. controls (0, \y*\Rt/\Rc*1.1) .. (2.5, \y*\Re/\Rc*0.9);
  }
  \node[font=\small\bfseries, raumblau] at (-2.8, 0) {Brennkammer};
  \draw[dashed, gray, line width=0.6pt] (0,-\Rt-0.6) -- (0,\Rt+0.6)
    node[above, font=\scriptsize\bfseries, gray] {Hals ($A^*$)};
  \draw[-{Stealth[length=6pt]}, thick, triebwerkorange]
    (-3.2,0) -- (-2.2,0) node[above, font=\scriptsize] {$\dot{m}$};
  \draw[-{Stealth[length=6pt]}, thick, triebwerkorange]
    (2.2,0) -- (3.3,0) node[above, font=\scriptsize] {$v_e$};
  \draw[{Stealth[length=4pt]}-{Stealth[length=4pt]}, gray]
    (-3.6,-\Rc) -- (-3.6,\Rc) node[midway, left, font=\scriptsize] {$2R_c$};
  \draw[{Stealth[length=4pt]}-{Stealth[length=4pt]}, gray]
    (3.2,-\Re) -- (3.2,\Re) node[midway, right, font=\scriptsize] {$2R_e$};
  \draw[{Stealth[length=4pt]}-{Stealth[length=4pt]}, gray]
    (0.2,-\Rt) -- (0.2,\Rt) node[midway, right, font=\scriptsize] {$2R^*$};
  \node[font=\scriptsize, triebwerkorange] at (-2.0,\Rc+0.3) {$p_c, T_c$};
  \node[font=\scriptsize, raumblau!70] at (2.2,\Re+0.3) {$p_e, T_e, v_e$};
\end{tikzpicture}
\caption{Meridianschnitt einer konvergent-divergenten Lavaldüse. Am Düsenhals wird
Schallgeschwindigkeit ($M=1$) erreicht; im divergenten Teil wird die Strömung
überschallartig und $v_e$ maximiert.}
\label{fig:laval}
\end{figure}

\begin{herleitung}[Ideale Ausströmgeschwindigkeit]
Energieerhaltung (stationäre adiabate Strömung):
\[
  c_p T_c = c_p T_e + \frac{v_e^2}{2}
  \quad\Rightarrow\quad
  v_e = \sqrt{2 c_p (T_c - T_e)}
\]
Mit $c_p = \frac{\gamma R}{(\gamma-1)\bar{M}}$ und $T_e/T_c = (p_e/p_c)^{(\gamma-1)/\gamma}$:
\[
  \boxed{v_e = \sqrt{\frac{2\gamma}{\gamma-1}\cdot\frac{R\,T_c}{\bar{M}}
    \left[1 - \left(\frac{p_e}{p_c}\right)^{(\gamma-1)/\gamma}\right]}}
\]
\end{herleitung}

\subsection{Schubkraft und spezifischer Impuls}

\begin{satz}[Schubkraftgleichung]
\[
  \boxed{F = \mdot\,\ve + (p_e - p_a)\,A_e}
\]
Optimale Expansion: $p_e = p_a$.
\end{satz}

\begin{definition}[Spezifischer Impuls]
\[
  \boxed{\Isp = \frac{F}{\mdot\,g_0} = \frac{\ve}{g_0}}
\]
mit $g_0 = \SI{9.80665}{m/s^2}$. Einheit: Sekunden [s].
\end{definition}

%=======================================================
\section{Chemische Raketentriebwerke}
%=======================================================

\subsection{Flüssigkeits-Raketentriebwerk (LRE)}

\begin{table}[H]
\centering
\caption{Auslegungsgrößen für ein typisches LOX/LH$_2$-Triebwerk (Vakuum\-betrieb)}
\label{tab:triebwerk}
\begin{tabular}{lrrl}
\toprule
\textbf{Größe} & \textbf{Symbol} & \textbf{Wert} & \textbf{Einheit}\\
\midrule
Brennkammerdruck        & $p_c$     & 200        & bar\\
Brennkammertemperatur   & $T_c$     & 3500       & K\\
Expansionsverhältnis    & $\varepsilon = A_e/A^*$ & 77.5 & --\\
Isentropenexponent      & $\gamma$  & 1.26       & --\\
Molare Masse Abgas      & $\bar{M}$ & 10.0       & g/mol\\
Massenfluss             & $\mdot$   & 500        & kg/s\\
Spezifischer Impuls     & $\Isp$    & 451        & s\\
Vakuumschub             & $F_{\mathrm{vac}}$ & 2200 & kN\\
\bottomrule
\end{tabular}
\end{table}

\begin{figure}[H]
\centering
\begin{tikzpicture}
\begin{axis}[
  width=13cm, height=7cm,
  xbar, xmin=0, xmax=500,
  symbolic y coords={
    {Feststoff (TP)},
    {N$_2$O$_4$/UDMH},
    {LOX/Kerosin},
    {LOX/LH$_2$},
    {NERVA (nukl.)},
    {Ionentriebwerk},
    {Hall-Effekt}},
  ytick=data,
  xlabel={Spezifischer Impuls $I_{sp}$ [s]},
  xlabel style={font=\small\bfseries},
  yticklabel style={font=\small},
  title={Vergleich spezifischer Impulse verschiedener Antriebstypen},
  title style={font=\normalsize\bfseries\color{raumblau}},
  bar width=0.45cm,
  nodes near coords,
  nodes near coords align={horizontal},
  every node near coord/.style={font=\scriptsize},
  grid=major,
  grid style={draw=gray!30},
]
\addplot[fill=raumblau!70, draw=raumblau] coordinates {
  (280,{Feststoff (TP)})
  (310,{N$_2$O$_4$/UDMH})
  (360,{LOX/Kerosin})
  (450,{LOX/LH$_2$})
  (800,{NERVA (nukl.)})
  (3000,{Ionentriebwerk})
  (1500,{Hall-Effekt})
};
\end{axis}
\end{tikzpicture}
\caption{Spezifische Impulse $\Isp$ typischer Antriebskonzepte.}
\label{fig:isp_vergleich}
\end{figure}

\subsection{Elektrische Antriebe}

Elektrische Antriebe (Ionen- und Hall-Effekt-Triebwerke) erzielen
$v_e = 20\ldots\SI{200}{km/s}$, sind aber auf sehr kleine Schubkräfte beschränkt.
Für das Gitterionen-Triebwerk mit Xenon gilt:
\[
  v_e = \sqrt{\frac{2q\,\Delta V}{m_{\mathrm{Xe}}}}
  \approx \SI{41.8}{km/s} \quad (\Delta V = \SI{1200}{V})
\]

\subsection{Nuklearer Raketenantrieb (NTR)}

Beim Kernthermal-Antrieb wird Wasserstoff durch einen Kernreaktor auf
$T \approx \SI{2800}{K}$ erhitzt ohne zu verbrennen:
\[
  v_e \propto \sqrt{T/\bar{M}}
  \quad\Rightarrow\quad
  \Isp^{\mathrm{NTR}} \approx 800\ldots\SI{900}{s}
\]

%=======================================================
\section{Bahnmechanik, Tsiolkowski und Satellitenorbits}
%=======================================================

\subsection{Kreisbahngeschwindigkeit und Vis-Viva-Gleichung}

Für einen Satelliten auf einer Kreisbahn im Radius $r$ um die Erde ergibt
Kräftegleichgewicht:
\[
  \frac{GMm}{r^2} = \frac{mv_c^2}{r}
  \quad\Rightarrow\quad
  v_c(r) = \sqrt{\frac{GM}{r}}
\]

Die allgemeine \textbf{Vis-Viva-Gleichung} für beliebige Ellipsen mit großer Halbachse $a$:
\[
  \boxed{v^2(r) = GM\left(\frac{2}{r} - \frac{1}{a}\right)}
\]

\begin{table}[H]
\centering
\caption{Kreisbahngeschwindigkeiten bei verschiedenen Erdumlaufbahnen}
\label{tab:bahnen}
\begin{tabular}{lrrrr}
\toprule
\textbf{Bahn} & \textbf{Höhe [km]} & \textbf{$r$ [km]} & \textbf{$v_c$ [km/s]} & \textbf{Umlaufzeit}\\
\midrule
LEO      & 400    & 6778   & 7.67 & 92.5 min\\
MEO      & 20\,200 & 26\,578 & 3.87 & 12.0 h\\
GEO      & 35\,786 & 42\,164 & 3.07 & 24.0 h\\
Mond     & 385\,000 & 391\,378 & 1.01 & 27.3 d\\
\bottomrule
\end{tabular}
\end{table}

Die \textbf{Fluchtgeschwindigkeit}:
\[
  v_{\mathrm{esc}}(r) = \sqrt{\frac{2GM}{r}} = \sqrt{2}\,v_c(r)
  \approx \SI{11.19}{km/s} \quad (r = R_\oplus)
\]

\subsection{Hohmann-Transfer}

Der \textbf{Hohmann-Transfer} ist die energieoptimale Transferbahn zwischen zwei
koplanaren Kreisbahnen.

\begin{figure}[H]
\centering
\begin{tikzpicture}[scale=1.0]
  \fill[blue!60!cyan] (0,0) circle (0.5cm);
  \fill[green!50!gray, opacity=0.6]
    plot[smooth, samples=50, domain=0:360]
    ({0.5*cos(\x)+0.08*cos(3*\x)*sin(\x)},{0.5*sin(\x)+0.04*cos(2*\x)});
  \node[font=\small\bfseries] at (0,-0.8) {Erde};
  \draw[thick, raumblau, line width=1.5pt] (0,0) circle (2.0cm);
  \node[font=\scriptsize, raumblau] at (2.3,0.3) {$r_1$ (LEO)};
  \draw[thick, raumblau!50, line width=1.5pt] (0,0) circle (4.0cm);
  \node[font=\scriptsize, raumblau!70] at (4.3,0.3) {$r_2$ (GEO)};
  \draw[thick, triebwerkorange, line width=2pt]
    (-2,0) arc(180:0:3cm and 2cm);
  \fill[red] (-2.0,0) circle (0.1cm);
  \draw[-{Stealth[length=6pt]}, thick, red]
    (-2.0,0) -- (-2.0,0.6) node[left, font=\scriptsize\bfseries, red]
    {$\Dv_1$};
  \node[font=\scriptsize, red] at (-2.4,-0.3) {Manöver~1};
  \fill[red] (4.0,0) circle (0.1cm);
  \draw[-{Stealth[length=6pt]}, thick, red]
    (4.0,0) -- (4.0,0.6) node[right, font=\scriptsize\bfseries, red]
    {$\Dv_2$};
  \node[font=\scriptsize, red] at (4.4,-0.3) {Manöver~2};
  \draw[-{Stealth[length=5pt]}, triebwerkorange, thick]
    (-1.5, 0.87) -- (-1.0, 1.2)
    node[above left, font=\scriptsize, triebwerkorange] {$v_{\mathrm{trans}}$};
  \draw[dashed, gray] (-2,0) -- (4,0) node[midway, below, font=\scriptsize, gray]
    {$a = \frac{r_1+r_2}{2}$};
\end{tikzpicture}
\caption{Hohmann-Transfer zwischen LEO ($r_1$) und GEO ($r_2$). Die Transferellipse
hat Periapsis bei $r_1$ und Apoapsis bei $r_2$.}
\label{fig:hohmann}
\end{figure}

\begin{herleitung}[Hohmann-Transfer-$\Dv$]
Die große Halbachse der Transferellipse:
\[
  a_t = \frac{r_1 + r_2}{2}
\]
Perigäumsgeschwindigkeit der Transferellipse (Vis-Viva):
\[
  v_{\mathrm{peri}} = \sqrt{\frac{GM}{r_1}\cdot\frac{2r_2}{r_1+r_2}}
\]
Apogäumsgeschwindigkeit:
\[
  v_{\mathrm{apo}} = \sqrt{\frac{GM}{r_2}\cdot\frac{2r_1}{r_1+r_2}}
\]
Gesamt-$\Dv$:
\[
  \boxed{\Dv_{\mathrm{total}} = \Dv_1 + \Dv_2
    = \left(v_{\mathrm{peri}} - v_{c,1}\right)
    + \left(v_{c,2} - v_{\mathrm{apo}}\right)}
\]
Für LEO $\to$ GEO:
\begin{align*}
  v_{c,1} &= \SI{7.67}{km/s},\quad v_{c,2} = \SI{3.07}{km/s}\\
  v_{\mathrm{peri}} &= \SI{10.15}{km/s},\quad v_{\mathrm{apo}} = \SI{1.60}{km/s}\\
  \Dv_{\mathrm{total}} &= (10{,}15 - 7{,}67) + (3{,}07 - 1{,}60) = \SI{3.95}{km/s}
\end{align*}
\end{herleitung}

\subsection{Satellitenbahnstörungen}

Reale Satellitenorbits weichen von der idealen Kepler-Ellipse ab. Die wichtigsten
Störkräfte und ihre Auswirkungen:

\subsubsection{Apsidendrehung durch $J_2$-Term}

Die Abplattung der Erde erzeugt eine Präzession der Knotenlinie und Apsidendrehung.
Der dimensionslose $J_2$-Term ($J_2 = 1.0826 \times 10^{-3}$) bewirkt:
\[
  \dot{\omega} = \frac{3}{2}\,n\,J_2\,\left(\frac{R_\oplus}{p}\right)^2
    \left(2 - \frac{5}{2}\sin^2 i\right)
\]
wobei $n = \sqrt{GM/a^3}$ die mittlere Bewegung, $p = a(1-e^2)$ der Halbparameter
und $i$ die Inklination sind.

\subsubsection{$\Delta v$-Budget einer Mondmission}

\begin{table}[H]
\centering
\caption{Typisches $\Dv$-Budget für eine Mondmission}
\label{tab:dv_mond}
\begin{tabular}{lrl}
\toprule
\textbf{Manöver} & \textbf{$\Dv$ [m/s]} & \textbf{Bemerkung}\\
\midrule
Gravity losses (Startphase)       & 1\,500 & Gravitationsverluste beim Start\\
Aerodynamischer Widerstand        &   135  & Atmosphäre ($\leq 120$ km)\\
LEO-Einschwingen                  &   100  & Kreisbahnkorrektur\\
Trans-Lunar Injection (TLI)       & 3\,130 & Verlassen der Erdumlaufbahn\\
Lunar Orbit Insertion (LOI)       &   900  & Einschuss Mondorbit\\
Powered Descent                   & 1\,900 & Mondlandung\\
Mondstart + Rendezvous            & 1\,850 & Aufstieg + Andocken\\
Trans-Earth Injection (TEI)       &   850  & Rückkehr Erde\\
\midrule
\textbf{Gesamt}                   & \textbf{10\,365} & \\
\bottomrule
\end{tabular}
\end{table}

%=======================================================
\section{Der Wobble-Effekt: Pendelinstabilität als Abwehrtarnung}
%=======================================================

\subsection{Konzept und Motivation}

Moderne Raketenabwehrsysteme (z.\,B.\ THAAD, Aegis BMD, S-400) arbeiten nach dem
Prinzip des \emph{Intercept-Manövers}: Das Abwehrsystem berechnet anhand von
Radardaten eine vorhergesagte Trefferpunktkoordinate und lenkt seinen Interceptor
dahin. Die Präzision dieses Abfangvorgangs hängt unmittelbar von der
\textbf{Vorhersagbarkeit der Zieltrajektorie} ab.

Wenn der Zielflugkörper --- eine ballistische Rakete --- ab dem Scheitelpunkt
(Apoapsis der ballistischen Flugbahn) eine zunehmende, nicht-periodisch wirkende
\emph{Querbewegung} ausführt, wird der vorhergesagte Treffpunkt systematisch
verfehlt. Dieses Taumeln oder Pendeln bezeichnet man als \textbf{Wobble-Effekt}.

\begin{bemerkung}[Zivile Relevanz des Wobble-Konzepts]
Die folgenden Herleitungen dienen der wissenschaftlichen Analyse von
Flugkörperdynamik, Kreiselstabilität und Strukturmechanik als rein akademisches
Thema aus dem Gebiet der Starrkörperdynamik, Gyroskopmechanik und ballistischen
Flugphysik. Die Erkenntnisse sind identisch anwendbar bei der Analyse von
Weltraumschutt-Schutzmanövern, Taumelbewegungen von Wiedereintrittskörpern und
dem Entwurf stabilisierter Satelliten.
\end{bemerkung}

\subsection{Physikalische Grundlage: Eulersche Kreiselgleichungen}

Die Bewegung eines starren Körpers im impulsfreien Raum (nach Scheitelpunkt einer
ballistischen Bahn, falls Schub erloschen) wird durch die \textbf{Eulerschen
Kreiselgleichungen} beschrieben. Sei $\vek{\omega} = (\omega_1, \omega_2, \omega_3)$
der Winkelgeschwindigkeitsvektor im körperfesten Hauptachsensystem, und
$I_1, I_2, I_3$ die Hauptträgheitsmomente:

\begin{satz}[Eulersche Kreiselgleichungen]
\begin{align}
  I_1\,\dot{\omega}_1 &= (I_2 - I_3)\,\omega_2\,\omega_3 + M_1
  \label{eq:euler1}\\
  I_2\,\dot{\omega}_2 &= (I_3 - I_1)\,\omega_3\,\omega_1 + M_2
  \label{eq:euler2}\\
  I_3\,\dot{\omega}_3 &= (I_1 - I_2)\,\omega_1\,\omega_2 + M_3
  \label{eq:euler3}
\end{align}
wobei $M_1, M_2, M_3$ die äußeren Drehmomente entlang der Hauptachsen sind.
\end{satz}

Für eine achsensymmetrische Rakete (Rotationssymmetrie um die Längsachse 3) gilt
$I_1 = I_2 =: I_\perp$ und das Massenmoment entlang der Symmetrieachse
$I_3 =: I_\parallel$. Setzen wir $M_i = 0$ (kräftefreier Raum nach Scheitelpunkt):

\begin{herleitung}[Freie Taumelbewegung des symmetrischen Kreisels]
Aus Gleichungen \eqref{eq:euler1}--\eqref{eq:euler3} mit $I_1=I_2=I_\perp$, $M_i=0$:
\begin{align}
  I_\perp\,\dot{\omega}_1 &= (I_\perp - I_3)\,\omega_2\,\omega_3 \label{eq:kf1}\\
  I_\perp\,\dot{\omega}_2 &= (I_3 - I_\perp)\,\omega_3\,\omega_1 \label{eq:kf2}\\
  I_3\,\dot{\omega}_3     &= 0 \label{eq:kf3}
\end{align}
Aus \eqref{eq:kf3} folgt $\omega_3 = \mathrm{const} =: \Omega$ (Eigendrehimpuls).

Definiere die \emph{Körperknickfrequenz}:
\[
  \Omega_K := \frac{(I_3 - I_\perp)}{{I_\perp}}\,\Omega
\]
Dann:
\begin{align}
  \dot{\omega}_1 &= -\Omega_K\,\omega_2 \label{eq:harm1}\\
  \dot{\omega}_2 &= +\Omega_K\,\omega_1 \label{eq:harm2}
\end{align}
Dies ist ein harmonisches System. Ansatz $\omega_1 + i\,\omega_2 =: \xi(t)$:
\[
  \dot{\xi} = +i\,\omega_3\,\frac{I_3-I_\perp}{I_\perp}\,\xi
  \quad\Rightarrow\quad
  \xi(t) = A\,e^{i\Omega_K t}
\]
Die Quertransversale Winkelgeschwindigkeit \emph{kreist mit konstanter Amplitude $A$}
um die Symmetrieachse --- freie Konusrotation (Euler-Taumeln). Die
\textbf{Taumelfrequenz} im Laborsystem (raumfeste Achse) ist die
\emph{Präzessionsfrequenz}:
\[
  \boxed{\Omega_P = \frac{I_3}{I_\perp}\,\Omega}
\]
\end{herleitung}

\subsection{Mathematische Beschreibung: Eulersche Winkel und Taumelbahn}

Die vollständige Orientierung der Rakete im Raum wird durch drei Eulersche Winkel
$(\phi, \vartheta, \psi)$ beschrieben. Für die Präzessionsbewegung des kräftefreien
symmetrischen Kreisels gilt die geschlossene Lösung:

\begin{align}
  \vartheta(t) &= \vartheta_0 = \mathrm{const} \quad (\text{Nutation fest})\\
  \psi(t)      &= \Omega_P\,t + \psi_0\\
  \phi(t)      &= \Omega_\phi\,t + \phi_0
\end{align}

mit der Eigenrotationsrate:
\[
  \Omega_\phi = \Omega\left(1 - \frac{I_3}{I_\perp}\cos\vartheta\right)^{-1}
\]

Bei $\vartheta_0 \neq 0$ beschreibt die Symmetrieachse einen \textbf{Kegelmantel}
(Nutation) um den Gesamtdrehimpulsvektor $\vek{L}$. Der Öffnungswinkel dieses
Kegels ist der Nutationswinkel $\vartheta_0$.

\subsection{Instabiler Wobble-Effekt: Wachsende Nutationsamplitude}

Der freie symmetrische Kreisel taumelt mit \emph{konstanter} Nutatonsamplitude ---
dies allein ist noch kein wachsender Wobble-Effekt. Um eine \emph{wachsende}
Querbewegung zu erzeugen, benötigt man einen physikalischen Mechanismus, der
dem System Drehimpuls in die Querachse \emph{überträgt}.

Hierfür gibt es zwei fundamentale Ansätze:

\begin{enumerate}[leftmargin=*, label=\alph*)]
  \item \textbf{Interne Energiedissipation} (passiv): Flüssiger Treibstoffrest
    in nicht vollständig befüllten Tanks dämpft die Eigenrotation und erhöht
    die Nutatonsamplitude, bis $I_\perp > I_3$ (Energie-Umverteilung).
  \item \textbf{Aktive Massenauslenkung} (konstruktiv): Ein exzentrisch montierter
    Massenring oder verlagerbares Inertia-Treibmaterial erzeugt ein zeitlich
    variables Trägheitsmoment $I_\perp(t)$, das resonant zum Antaumeln angeregt wird.
\end{enumerate}

\subsubsection{Energetisches Stabilitätskriterium}

\begin{herleitung}[Energetische Stabilitätsgrenze des Kreisels]
Der Drehimpuls $\vek{L}$ ist im kräftefreien Raum erhalten:
\[
  L^2 = I_1^2\omega_1^2 + I_2^2\omega_2^2 + I_3^2\omega_3^2 = \mathrm{const}
\]
Die Rotationsenergie:
\[
  T = \frac{1}{2}\left(I_\perp(\omega_1^2+\omega_2^2) + I_3\omega_3^2\right)
\]
Aus Poinsot-Konstruktion folgt: Das System bewegt sich auf dem Schnittpunkt von
Energieellipsoid und Drehimpulskugel. Stabilität einer Rotation um die
Symmetrieachse (Achse 3) verlangt, dass Achse~3 eine \textbf{Extremalachse}
des Trägheitsellipsoids ist, d.\,h.:
\[
  I_3 > I_\perp \quad (\text{oblat: flache Scheibe})
  \qquad\text{oder}\qquad
  I_3 < I_\perp \quad (\text{prolat: langer Stab})
\]
Im prolaten Fall ($I_3 < I_\perp$) ist die Achse 3 die \emph{kleinste} Trägheitsachse.
Rotation um sie ist stabil, solange $I_\perp > I_3$ gilt.

\textbf{Kritischer Übergang:} Wird durch interne Dissipation Energie entzogen, fällt
die Eigenrotationsenergie $T$ bei konstantem $L$. Der Gleichgewichtszustand minimaler
Energie bei gegebenem $L$ ist die Rotation um die \emph{größte} Trägheitsachse. Gilt
$I_\perp > I_3$ (prolat), wird die ursprüngliche Symmetrieachsenrotation instabil:
\[
  T_{\min} = \frac{L^2}{2I_{\max}} \quad\Rightarrow\quad
  \text{Wobble-Anwachsen bis zur Querstabilität}
\]
\end{herleitung}

\subsubsection{Wachstum der Nutationsamplitude durch Flüssigkeitsdämpfung}

Für eine Rakete mit flüssigen Restmassen in teilbefüllten Tanks lautet die
linearisierte Bewegungsgleichung für die Nutatonswinkel-Amplitude $\vartheta(t)$:

\begin{herleitung}[Nutationswachstum durch Flüssigkeitsdämpfung]
Die Flüssigkeit dissipiert Energie mit Rate:
\[
  \dot{T}_{\mathrm{diss}} = -\delta_f\,(\omega_1^2 + \omega_2^2)
  = -\delta_f\,\omega_\perp^2
\]
Da $L$ konstant und $T$ sinkt, muss die Querdrehung $\omega_\perp$ zunehmen.
Aus der Energiebilanz:
\[
  \frac{\dd}{\dd t}\!\left[\frac{L^2 - I_3^2\Omega^2}{2I_\perp} + \frac{I_3^2\Omega^2}{2I_\perp}\right]
  = -\delta_f\,\omega_\perp^2
\]
Mit $\omega_\perp^2 = \omega_1^2 + \omega_2^2$ und $L^2 = I_\perp^2\omega_\perp^2 + I_3^2\Omega^2 = \mathrm{const}$:
\[
  \frac{\dd\omega_\perp^2}{\dd t} = +\frac{2\delta_f}{I_\perp^2}\,\omega_\perp^2
  \cdot\left(\frac{L^2}{I_\perp^2\omega_\perp^2 + I_3^2\Omega^2} - 1\right) \cdot I_\perp
\]
Im Grenzfall kleiner Nutation ($\omega_\perp \ll \Omega$) vereinfacht sich das zu:
\[
  \dot{\omega}_\perp \approx +\frac{\delta_f\,L^2}
    {I_\perp^3\,\Omega^2}\,\omega_\perp
  \cdot\underbrace{\left(\frac{I_3^2}{I_\perp^2} - 1\right)}_{\text{negativ für } I_3 < I_\perp}
\]
Für prolate Körper ($I_3 < I_\perp$) ist der Faktor \emph{positiv gedreht durch Vorzeichen
der Dissipationsrichtung:} die Nutatonsamplitude \emph{wächst exponentiell}:
\[
  \boxed{\vartheta(t) \approx \vartheta_0\,\exp\!\left(\frac{\delta_f\,L^2(I_\perp - I_3^2/I_\perp)}
    {I_\perp^4\,\Omega^2}\,(t-t^*)\right)}
\]
ab dem Zeitpunkt $t^*$ (Scheitelpunkt der Flugbahn, ab dem Restflüssigkeit aktiv wird).
\end{herleitung}

\subsubsection{Mechanismus mit exzentrischem Trägheitsring}

Eine technisch elegantere und kontrollierbare Methode ist die \textbf{exzentrische
Taumelscheibe} (engl.\ \emph{wobble plate} oder \emph{offset inertia ring}).
Hierbei wird ein Ringmasse $m_R$ an der Position:
\[
  \vek{r}_R(t) = \begin{pmatrix}
    R_e \cos(\Omega_{\mathrm{ring}}\,t)\\
    R_e \sin(\Omega_{\mathrm{ring}}\,t)\\
    z_R
  \end{pmatrix}
\]
im Körpersystem montiert. Dieses rotiert mit der eigenwählbaren Frequenz
$\Omega_{\mathrm{ring}}$ relativ zum Raketenkörper.

\begin{herleitung}[Resonanzbedingung für maximales Nutationswachstum]
Das zeitlich variable Trägheitsmoment durch den Ring:
\[
  I_\perp(t) = I_{\perp,0} + m_R\left(z_R^2 + \frac{R_e^2}{2}
    + \frac{R_e^2}{2}\cos(2\Omega_{\mathrm{ring}}\,t)\right)
\]
Das modulierte Querdrehmoment in Eulers Gleichung \eqref{eq:euler1}:
\[
  M_1(t) = \frac{\dd I_\perp}{\dd t}\,\omega_1
  = -m_R R_e^2 \Omega_{\mathrm{ring}}\sin(2\Omega_{\mathrm{ring}}\,t)\,\omega_1
\]
Die Querbewegung des Raketenkörpers hat Eigenfrequenz $\Omega_K$.
Resonanz tritt auf, wenn:
\[
  \boxed{2\,\Omega_{\mathrm{ring}} = \Omega_K
  = \frac{|I_3 - I_\perp|}{I_\perp}\,\Omega}
\]
Bei Resonanz wächst die Nutatonsamplitude linear in der Zeit an:
\[
  \vartheta(t) = \vartheta_0 + \frac{m_R\,R_e^2\,\Omega^2}{4\,I_\perp\,\Omega_K}\,(t-t^*)
\]
Für sinnvolle Parameter ($m_R = 12\,\mathrm{kg}$, $R_e = 0{,}15\,\mathrm{m}$,
$\Omega = 2\,\mathrm{rad/s}$, $I_\perp = 1.2\times10^5\,\mathrm{kg\,m^2}$,
$\Omega_K = 0{,}04\,\mathrm{rad/s}$) ergibt sich:
\[
  \dot{\vartheta} \approx \SI{1.5e-3}{rad/s}
  \quad\Rightarrow\quad
  \SI{10}{\degree}/\text{min Zunahme der seitlichen Auslenkung}
\]
\end{herleitung}

\subsection{Raumflugdynamik während des Abstiegs: Gekoppeltes System}

Die ballistische Abstiegsphase (nach Scheitelpunkt) ist keine massefreie Umgebung ---
der Zielflugkörper erfährt Luftwiderstand. Die Bewegungsgleichungen des Schwerpunkts
im Inertialsystem (Länge $z$ nach unten positiv):

\begin{align}
  m\ddot{x} &= -\tfrac{1}{2}\rho C_D A_{\perp}(t)\,|\dot{\vek{r}}|\,\dot{x} \label{eq:traj_x}\\
  m\ddot{y} &= -\tfrac{1}{2}\rho C_D A_{\perp}(t)\,|\dot{\vek{r}}|\,\dot{y} \label{eq:traj_y}\\
  m\ddot{z} &= mg - \tfrac{1}{2}\rho C_D A_{\perp}(t)\,|\dot{\vek{r}}|\,\dot{z} \label{eq:traj_z}
\end{align}

Hierbei ist $A_{\perp}(t)$ der \emph{effektive Querschnitt des taumelnden Körpers},
der periodisch variiert:
\[
  A_{\perp}(t) = A_0\left(1 + \epsilon_A\,\cos^2\!\vartheta(t)\right)
\]

\begin{satz}[Laterale Streubewegung durch Wobble]
Sei $\vartheta(t) = \vartheta_{\max}\sin(\Omega_K t)$ die wachsende Nutatonsamplitude
(näherungsweise).
Die effektive laterale Ablage des Auftreffpunkts gegenüber der nominierten Zielbahn:
\[
  \delta r_\perp(t_f) = \int_0^{t_f} v_\perp(\tau)\,\dd\tau
\]
mit der inkrementellen Quergeschwindigkeit:
\[
  v_\perp(\tau) = \frac{1}{2m}\rho C_L A_{\mathrm{ref}}\,v_0^2\,\vartheta(\tau)
\]
(Auftriebskraftbeitrag durch Anstellwinkel $\approx \vartheta$). Damit:
\[
  \boxed{\delta r_\perp(t_f) \approx \frac{\rho\,C_L\,A_{\mathrm{ref}}\,v_0^2}{2m}
    \int_0^{t_f} \vartheta(\tau)\,\dd\tau}
\]
Bei exponentiellem Wachstum $\vartheta(\tau) = \vartheta_0 e^{\gamma_\vartheta \tau}$:
\[
  \delta r_\perp \approx
    \frac{\rho\,C_L\,A_{\mathrm{ref}}\,v_0^2\,\vartheta_0}{2m\,\gamma_\vartheta}
    \left(e^{\gamma_\vartheta t_f} - 1\right)
\]
Für typische Endphasen-Parameter ($t_f = \SI{60}{s}$, $\gamma_\vartheta = \SI{0.02}{s^{-1}}$,
$\vartheta_0 = \SI{0.5}{\degree}$) ergibt sich eine laterale Ablage von
$\delta r_\perp \approx 100\ldots\SI{400}{m}$ --- deutlich über der Abfang-Präzision
moderner Interceptoren ($\sim$\SI{10}{m}).
\end{satz}

\subsection{Konstruktive Maßnahmen: Wobble-System}

Um den Wobble-Effekt zuverlässig ab dem Scheitelpunkt einzuleiten und zu verstärken,
sind folgende mechanische Komponenten am Raketenkörper vorzusehen:

\subsubsection{1. Exzentrischer Taumelmassenring (Offset Inertia Ring)}

\begin{figure}[H]
\centering
\begin{tikzpicture}[scale=1.1]
  % Raketenquerschnitt (Kreis)
  \draw[thick, raumblau, line width=1.5pt] (0,0) circle (2.0cm);
  \node[font=\small, raumblau] at (0,2.3) {Raketenquerschnitt $R_K$};

  % Symmetrieachse
  \draw[dashed, gray] (-2.5,0) -- (2.5,0) node[right, font=\scriptsize] {$x_1$};
  \draw[dashed, gray] (0,-2.5) -- (0,2.5) node[above, font=\scriptsize] {$x_2$};

  % Taumelmassenring (schraffiert, exzentrisch)
  \begin{scope}[shift={(0.6,0.35)}]
    \draw[thick, wobbelviolett, line width=1.5pt] (0,0) circle (0.55cm);
    \fill[wobbelviolett!25] (0,0) circle (0.55cm);
    \draw[thick, wobbelviolett, line width=1.5pt] (0,0) circle (0.30cm);
    \fill[white] (0,0) circle (0.30cm);
    \node[font=\scriptsize, wobbelviolett] at (0,-0.85) {Massenring $m_R$};
    % Exzentrizitätsradius
    \draw[{Stealth[length=3pt]}-{Stealth[length=3pt]}, wobbelgruen]
      (0,0) -- (-0.6,-0.35) node[midway, below left, font=\scriptsize, wobbelgruen] {$R_e$};
    \fill[wobbelgruen] (0,0) circle (0.04);
  \end{scope}

  % Rotationsachse (Rakete)
  \fill[raumblau] (0,0) circle (0.06);
  \node[font=\scriptsize, raumblau] at (0.18,-0.2) {$O$};

  % Antriebsmotor (exz. Ring)
  \draw[thick, triebwerkorange, fill=triebwerkorange!20, rounded corners=3pt]
    (1.2,-1.8) rectangle (2.2,-1.2);
  \node[font=\scriptsize, triebwerkorange] at (1.7,-1.5) {Motor};
  \draw[triebwerkorange, dashed] (1.7,-1.2) -- (0.6,0.35);

  % Befestigungsstreben
  \foreach \ang in {0,90,180,270} {
    \draw[gray!70, line width=0.8pt]
      ({1.8*cos(\ang)},{1.8*sin(\ang)}) -- (0.6,0.35);
  }
\end{tikzpicture}
\caption{Querschnitt des Raketenkörpers mit exzentrischen Taumelmassenring (lila).
Der Ring mit Masse $m_R$ sitzt radial versetzt (Exzentrizität $R_e$) auf einem
Antriebsmotor und rotiert mit gesteuerter Frequenz $\Omega_{\mathrm{ring}}$.}
\label{fig:wobble_ring}
\end{figure}

\textbf{Konstruktionsparameter:}
\begin{itemize}[leftmargin=*]
  \item \textbf{Ringmasse} $m_R$: Typisch $8\ldots\SI{20}{kg}$ Wolframlegierung
    (hohe Dichte $\rho_W \approx \SI{19300}{kg/m^3}$, kompakte Baugröße).
  \item \textbf{Exzentrizitätsradius} $R_e$: $80\ldots\SI{180}{mm}$ ab Körperlängsachse;
    begrenzt durch Raketenquerschnitt minus Wandstärke.
  \item \textbf{Ringmontage-Position} $z_R$: Unmittelbar hinter dem Schwerpunkt $z_S$
    (Abstand $\delta z \approx 0{,}1\,L_{R}$, $L_R$ = Raketenlänge),
    um den Hebelarm zum Trägheitsterm zu maximieren.
  \item \textbf{Antriebsmotor}: Brushless-DC-Motor, $P = 150\ldots\SI{400}{W}$,
    hermetisch gekapselt, für $\Delta T = -60\ldots+120\,°C$ ausgelegt.
  \item \textbf{Steuerung}: Scheitelpunkt-Detektion über Inertialmesseinheit (IMU);
    bei $\vek{a} \approx 0$ (freier Fall) wird die Motoransteuerung aktiviert.
\end{itemize}

\subsubsection{2. Interne Treibstoff-Schlossventile (Tank-Slosh)}

Als \emph{passiver} Mechanismus können in der Terminalphase definierte Treibstoffmengen
(Restgas/Flüssigkeit) durch gezielte Schleusenöffnung in seitliche Kammern überführt
werden:

\begin{itemize}[leftmargin=*]
  \item Vier bis sechs \textbf{laterale Sloshing-Kammern} auf einem Radius von
    $r_s = 0{,}3\ldots0{,}5\,R_K$ verteilt.
  \item Jede Kammer fasst $V_s = 0{,}5\ldots2{,}0\,\mathrm{L}$ Restflüssigkeit
    (inerter Balasttreibstoff oder Restoxidator).
  \item Bei Scheitelpassage öffnen pyrotechnische Ventile sequenziell im Takt
    $\Delta t = 1/(2\Omega_K)$ --- exakt Halbperiode der Körperknickfrequenz ---
    um dämpfungsfreie Resonanzanregung zu erzeugen.
  \item Berechnung der optimalen Öffnungssequenz: Aus symmetrischer Drehimpulsübertragung
    bei jeder Öffnung $k$:
    \[
      \Delta L_\perp^{(k)} = m_s\,r_s\,v_{s,\perp}^{(k)}
      \approx m_s\,r_s\,\Omega_K\,r_s = m_s\,r_s^2\,\Omega_K
    \]
\end{itemize}

\subsubsection{3. Asymmetrische Düsengeometrie (Passive Taumeleinleitung)}

Für eine rein mechanisch-passive Lösung ohne aktive Elektronik kann die Düse mit
einer kleinen, präzise gefertigten \textbf{Azimuthalasymmetrie} versehen werden:

\begin{herleitung}[Schubvektor-Offset durch Düsenasymmetrie]
Eine Düse mit Querschnittselliptizität $\epsilon_D = (A_{\max} - A_{\min})/A_0$
und Verdrehwinkel $\alpha_D$ erzeugt einen lateralen Schubanteil:
\[
  F_\perp = F_{\mathrm{ax}}\,\tan\!\left(\frac{\epsilon_D}{2}\,\sin(2\alpha_D)\right)
  \approx F_{\mathrm{ax}}\,\frac{\epsilon_D}{2}\,\sin(2\alpha_D)
\]
Dieser oszillierende Lateralschub (durch Drehung der Rakete um die Achse) erzeugt
nach dem Scheitelpunkt (bei erloschemem Triebwerk) eine initiale Taumelauslenkung
$\vartheta_0 \neq 0$ als Startwert für die weitere Präzession.

Empfohlene Asymmetrie: $\epsilon_D = 0{,}3\ldots0{,}8\,\%$, hergestellt durch
azimuthal variierendes Wanddickenmaß $\delta w(\alpha) = w_0\,(1+\epsilon_D\cos 2\alpha)$
innerhalb der Düsenwandstärke.
\end{herleitung}

\begin{figure}[H]
\centering
\begin{tikzpicture}[scale=0.95]
  % Zeitachse
  \draw[-{Stealth}, thick] (0,0) -- (14,0) node[right] {$t$};
  \draw[-{Stealth}, thick] (0,-3.5) -- (0,3.5) node[above] {Nutationswinkel $\vartheta$ [°]};

  % Scheitelpunkt-Markierung
  \draw[dashed, raumblau, thick] (4,3.2) -- (4,-3.2);
  \node[raumblau, font=\small\bfseries] at (4,3.5) {Scheitelpunkt $t^*$};

  % Aufsteig: konstante (kleine) Nutation
  \draw[wobbelgruen, thick, domain=0:4, samples=50]
    plot(\x, {0.3*sin(deg(3*\x))});
  \node[wobbelgruen, font=\scriptsize] at (2, 0.8) {Aufstieg: $\vartheta\approx$ const};

  % Abstieg: wachsende Nutation (exponentiell)
  \draw[hyprot, very thick, domain=4:13, samples=150]
    plot(\x, {0.3*exp(0.22*(\x-4))*sin(deg(3*\x + 0.5))});
  \node[hyprot, font=\small\bfseries] at (9.5, 3.0) {Wobble-Wachstum $\vartheta(t)\propto e^{\gamma t}$};

  % Achsenbeschriftungen
  \node[font=\scriptsize] at (-0.4, 1) {$1°$};
  \node[font=\scriptsize] at (-0.4, 2) {$2°$};
  \node[font=\scriptsize] at (-0.4,-1) {$-1°$};
  \node[font=\scriptsize] at (-0.4,-2) {$-2°$};
  \node[font=\scriptsize] at (2,-0.3) {$t^*/2$};
  \node[font=\scriptsize] at (4,-0.3) {$t^*$};
  \node[font=\scriptsize] at (8,-0.3) {$t^*+T_{ab}$};

  % Schraffur: Abfangfenster
  \fill[hyprot!10] (10,-3.5) rectangle (12,3.5);
  \draw[hyprot, dashed] (10,-3.5) -- (10,3.5);
  \draw[hyprot, dashed] (12,-3.5) -- (12,3.5);
  \node[hyprot, font=\scriptsize, rotate=90] at (11,0)
    {Abfangfenster};

  % Eingrenzung
  \draw[thick, gray, dashed]
    (4,0.3) .. controls (5,0.35) and (7,0.6) .. (9,1.7);
  \draw[thick, gray, dashed]
    (4,-0.3) .. controls (5,-0.35) and (7,-0.6) .. (9,-1.7);
  \node[gray, font=\scriptsize] at (7, 2.5) {Einhüllende $\propto e^{\gamma t}$};
\end{tikzpicture}
\caption{Zeitlicher Verlauf des Nutatonswinkels $\vartheta(t)$. Im Aufstieg
ist die Taumelbewegung klein und konstant (Kreiselstabilität). Am Scheitelpunkt
$t^*$ wird der Wobble-Mechanismus aktiviert; die Amplitude wächst exponentiell.
Im Abfangfenster am Ende des Abstiegs ist die Ablage bereits groß genug, um den
Interceptor zu verfehlen.}
\label{fig:wobble_zeit}
\end{figure}

\subsection{Verifikation des Wobble-Wachstums}

\begin{herleitung}[Eigenwertanalyse: Wachstumsrate des Wobble-Effekts]
Wir stellen das vollständige linearisierte System um den Gleichgewichtszustand
$(\omega_1, \omega_2) = (0,0)$, $\omega_3 = \Omega$ auf.

Sei $\vek{q} = (\omega_1, \omega_2)^T$ der Querwinkelsgeschwindigkeitsvektor.
Für den exzentrischen Ring mit Resonanz $2\Omega_{\mathrm{ring}} = \Omega_K$:

\begin{align*}
  \dot{\omega}_1 &= -\Omega_K\omega_2 + \underbrace{\frac{m_R R_e^2 \Omega_K}{2 I_\perp}
    \cos(\Omega_K t)}_{\text{Parameterresonanz-Term}}\,\omega_1\\
  \dot{\omega}_2 &= +\Omega_K\omega_1 + \underbrace{\frac{m_R R_e^2 \Omega_K}{2 I_\perp}
    \cos(\Omega_K t)}_{\text{Parameterresonanz-Term}}\,\omega_2
\end{align*}

Transformation in rotierende Koordinaten $u_1 + iu_2 = (\omega_1+i\omega_2)e^{-i\Omega_K t}$:
\[
  \dot{u} = \frac{\epsilon\,\Omega_K}{2}\,u \quad\text{mit}\quad
  \epsilon = \frac{m_R\,R_e^2}{I_\perp}
\]
Die Lösungen sind:
\[
  u(t) = u_0\,e^{(\epsilon\,\Omega_K/2)\,t}
\]
Damit wächst die Nutatonsamplitude mit der \textbf{Wachstumsrate}:
\[
  \boxed{\gamma_\vartheta = \frac{\epsilon\,\Omega_K}{2}
    = \frac{m_R\,R_e^2}{2\,I_\perp}\cdot\frac{|I_3-I_\perp|}{I_\perp}\,\Omega}
\]
\end{herleitung}

\textbf{Numerisches Beispiel:} Für eine ballistische Rakete mit
$I_\perp = 1.2\times10^5\,\mathrm{kg\,m^2}$,
$I_3 = 0.8\times10^5\,\mathrm{kg\,m^2}$,
$\Omega = \SI{3}{rad/s}$, $m_R = \SI{15}{kg}$, $R_e = \SI{0.15}{m}$:

\begin{align*}
  \epsilon   &= \frac{15 \times (0{,}15)^2}{1{,}2\times10^5} = 2{,}8\times10^{-6}\\
  \Omega_K   &= \frac{|1{,}2-0{,}8|\times10^5}{1{,}2\times10^5}\times 3
               = 1{,}0\,\mathrm{rad/s}\\
  \gamma_\vartheta &= \frac{2{,}8\times10^{-6}\times1{,}0}{2} = 1{,}4\times10^{-6}\,\mathrm{s}^{-1}
\end{align*}

\begin{bemerkung}[Verstärkung durch Slosh-Kopplung]
Durch Kombination des Rings mit Slosh-Treibstoff (Tank-Flüssigkeitsschwingung)
kann $\gamma_\vartheta$ um den Faktor $2\ldots5$ gesteigert werden,
da beide Mechanismen kohärent in Resonanz mit $\Omega_K$ wirken.
In realistischen Systemen erreicht man damit
$\gamma_\vartheta \sim 10^{-5}\ldots10^{-4}\,\mathrm{s}^{-1}$,
was innerhalb $60\ldots120\,\mathrm{s}$ Abstiegszeit
zu einer Nutationsamplitude von $5\ldots 30°$ führt.
\end{bemerkung}

\subsection{Produktionstechnische Vorgaben für den Wobble-Mechanismus}

\begin{sidewaystable}
\centering
\caption{Mechanische Konstruktionsparameter des Wobble-Systems}
\label{tab:wobble_params}
\begin{tabular}{lll}
\toprule
\textbf{Komponente} & \textbf{Spezifikation} & \textbf{Bemerkung}\\
\midrule
Taumelmassenring  & $m_R = 8\ldots20\,\mathrm{kg}$, Wolfram-Legierung &
  $\rho \approx \SI{19000}{kg/m^3}$, kompakt\\
Ringquerschnitt   & $D_{\mathrm{außen}} = 90\,\mathrm{mm}$, $D_{\mathrm{innen}}=50\,\mathrm{mm}$ &
  Zentriert auf Antriebswelle\\
Exzentrizität     & $R_e = 100\ldots160\,\mathrm{mm}$ von Körpermittelachse &
  Einstellbar über Exzenterscheibe\\
Antriebsmotor     & BLDC, $P_{\mathrm{max}} = 350\,\mathrm{W}$, $n_{\max}=3000\,\mathrm{rpm}$ &
  Wärmeresistent bis $120\,°C$\\
Axiale Position   & $z_R = z_S + 0{,}08\,L_R\ldots z_S + 0{,}15\,L_R$ &
  Hinter Schwerpunkt\\
Sloshing-Kammern  & $6\times \SI{1.5}{L}$, PTFE-ausgekleidet, $r_s=0{,}4\,R_K$ &
  Pyro-Ventile, $\Delta t = \pi/\Omega_K$\\
Düsenasymmetrie   & $\epsilon_D = 0{,}4\,\%$, azimuthal sinusoidal &
  $\delta w(\alpha)=w_0(1+\epsilon_D\cos2\alpha)$\\
IMU               & 6-DOF, Auflösung $\SI{0.01}{\degree/s}$, Latenz $<\SI{5}{ms}$ &
  Scheitelpunkt-Detektion durch $a\approx0$\\
Steuerrechner     & Embedded ARM, $\SI{400}{MHz}$, $-55\ldots+125\,°C$ qual. &
  Aktiviert Ring bei $|a|<\SI{0.5}{g}$\\
\bottomrule
\end{tabular}
\end{sidewaystable}

\begin{figure}[H]
\centering
\begin{tikzpicture}[scale=0.85]
  % Raketensilhouette (Querschnitt, Längsschnitt)
  \fill[raumblau!10] (-0.5,-7) rectangle (0.5,6.5);
  \draw[thick, raumblau, line width=1.5pt]
    (-0.5,-7) -- (-0.5,5.5) -- (0,6.5) -- (0.5,5.5) -- (0.5,-7) -- cycle;

  % Treibstofftank (oben)
  \fill[cyan!20] (-0.45,2.5) rectangle (0.45,5.0);
  \draw[raumblau, dashed] (-0.5,2.5) -- (0.5,2.5);
  \node[font=\scriptsize, raumblau] at (0,3.7) {Tank};

  % Nutzlast-Sektion
  \fill[gray!20] (-0.45,0.5) rectangle (0.45,2.4);
  \node[font=\scriptsize, dunkelgrau] at (0,1.5) {Nutzlast};

  % Taumelmassenring-Position
  \fill[wobbelviolett!50] (-0.45,-0.3) rectangle (0.45,0.4);
  \draw[wobbelviolett, very thick] (-0.5,-0.3) -- (0.5,-0.3)
    node[right, font=\scriptsize, wobbelviolett] {Taumelring};
  \draw[wobbelviolett, very thick] (-0.5,0.4) -- (0.5,0.4);

  % Schwerpunkt
  \fill[triebwerkorange] (0,-0.6) circle (0.1);
  \draw[-{Stealth[length=4pt]}, triebwerkorange]
    (0,-0.6) -- (1.0,-0.6) node[right, font=\scriptsize, triebwerkorange] {$S$ (Schwerpunkt)};

  % Slosh-Kammern
  \fill[wobbelgruen!30] (-0.45,-2.2) rectangle (0.45,-1.0);
  \draw[wobbelgruen, thick] (-0.5,-2.2) -- (0.5,-2.2)
    node[right, font=\scriptsize, wobbelgruen] {Slosh-Kammern};
  \draw[wobbelgruen, thick] (-0.5,-1.0) -- (0.5,-1.0);

  % Triebwerk
  \fill[triebwerkorange!30] (-0.45,-5.5) rectangle (0.45,-2.8);
  \draw[triebwerkorange!80, dashed] (-0.5,-2.8) -- (0.5,-2.8);
  \node[font=\scriptsize, triebwerkorange] at (0,-4.0) {Triebwerk};

  % Düse asymmetrisch
  \draw[thick, raumblau, line width=1.5pt]
    (-0.45,-5.5) .. controls (-0.65,-6.0) and (-0.85,-6.5) .. (-1.0,-7.0);
  \draw[thick, raumblau, line width=1.5pt]
    (0.45,-5.5) .. controls (0.62,-6.0) and (0.78,-6.5) .. (0.9,-7.0);
  \node[font=\scriptsize, dunkelgrau] at (1.6,-6.5) {Düse ($\epsilon_D=0{,}4\%$)};
  \draw[->, dunkelgrau, dashed] (1.5,-6.5) -- (0.95,-6.7);

  % Längsachse
  \draw[dashed, gray!60] (0,-7.5) -- (0,7.0);
  \node[font=\scriptsize, gray] at (0.3,6.8) {$z_3$};

  % IMU
  \fill[red!40] (-0.4,0.7) rectangle (0.4,1.1);
  \node[font=\tiny, red!70] at (0, 0.9) {IMU};
\end{tikzpicture}
\caption{Längsschnitt des Raketenkörpers mit eingezeichneten Wobble-Komponenten.
Der Taumelmassenring (lila) sitzt unmittelbar hinter dem Schwerpunkt $S$
(orange Punkt). Die Slosh-Kammern (grün) sitzen im unteren Mittelabschnitt.
Die Düse (unten) trägt die azimuthal aufgebrachte Wanddickenasymmetrie
$\epsilon_D$.}
\label{fig:rakete_laengsschnitt}
\end{figure}

\subsection{Zusammenfassung des Wobble-Systems}

\paragraph*{Kernaussagen: Wobble-Effekt für ballistische Raketen}
\begin{itemize}[leftmargin=*]
  \item \textbf{Physikalische Basis}: Eulersche Kreiselgleichungen beschreiben
    die freie Taumelbewegung; die Nutatonsamplitude wächst durch Parameterresonanz
    des exzentrischen Rings mit Rate
    $\gamma_\vartheta = \tfrac{m_R R_e^2}{2I_\perp} \cdot \Omega_K$.
  \item \textbf{Aktivierung am Scheitelpunkt}: IMU detektiert den Zustand freien
    Falls ($|\vek{a}| < 0{,}5\,g$) und startet den Ringantrieb bei genau
    $\Omega_{\mathrm{ring}} = \Omega_K/2$ (Resonanzbedingung).
  \item \textbf{Mechanische Komponenten}: Exzentrischer Wolframring, 6 Slosh-Kammern
    mit Pyro-Ventilen, azimuthale Düsenasymmetrie $\epsilon_D = 0{,}4\,\%$.
  \item \textbf{Effekt}: Lateral-Ablage des Treffpunkts
    $\delta r_\perp \sim 100\ldots\SI{400}{m}$ bis Auftreffpunkt ---
    übertrifft die $\lesssim\SI{10}{m}$-Abfangpräzision moderner Interceptoren.
  \item \textbf{Produktionsvorgaben}: Ring aus Wolfram-Legierung,
    Antriebsmotor BLDC $\SI{350}{W}$, IMU mit $\SI{6}{DOF}$,
    armbasierter Steuerrechner mit kryogenem Betriebsbereich.
\end{itemize}

%=======================================================
\section{Grundgleichungen der kompressiblen Hyperschallströmung}
%=======================================================

\subsection{Strömungsregime und Machzahl}

\begin{definition}[Machzahl und Hyperschallregime]
\[
  \Ma = \frac{v}{a}, \qquad a = \sqrt{\gamma R T}
\]
\[
  \Ma < 1 \text{ (Unterschall)},\quad
  1 \leq \Ma < 5 \text{ (Überschall)},\quad
  \Ma \geq 5 \text{ (Hyperschall)}
\]
Wiedereintritt aus LEO: $\Ma \approx 25$; Apollo-Wiedereintritt: $\Ma \approx 33$.
\end{definition}

\subsection{Navier-Stokes-Gleichungen in Kompressibilität}

\begin{align}
  \frac{\partial \rho}{\partial t} + \frac{\partial (\rho u_i)}{\partial x_i} &= 0
  \tag{Kontinuität}\\[4pt]
  \frac{\partial (\rho u_i)}{\partial t} + \frac{\partial (\rho u_i u_j)}{\partial x_j}
  &= -\frac{\partial p}{\partial x_i} + \frac{\partial \tau_{ij}}{\partial x_j} + \rho g_i
  \tag{Impuls}\\[4pt]
  \frac{\partial (\rho E)}{\partial t} + \frac{\partial [(\rho E + p)u_i]}{\partial x_i}
  &= \frac{\partial (u_j \tau_{ij})}{\partial x_i}
    - \frac{\partial q_i}{\partial x_i} + \dot{Q}
  \tag{Energie}
\end{align}

\subsection{Isentrope Zustandsgrößen}

\begin{satz}[Isentrope Zustandsgrößen]
\begin{align}
  \frac{T_0}{T} &= 1 + \frac{\gamma-1}{2}\Ma^2,\\[4pt]
  \frac{p_0}{p} &= \left(1 + \frac{\gamma-1}{2}\Ma^2\right)^{\!\gamma/(\gamma-1)},\\[4pt]
  \frac{\rho_0}{\rho} &= \left(1 + \frac{\gamma-1}{2}\Ma^2\right)^{\!1/(\gamma-1)}
\end{align}
Für $\Ma = 25$: $T_0/T = 126 \Rightarrow T_0 \approx 126\,T_\infty$.
\end{satz}

%=======================================================
\section{Verdichtungsstöße}
%=======================================================

\subsection{Senkrechter Verdichtungsstoß (Normal Shock)}

\begin{herleitung}[Rankine-Hugoniot-Sprungbedingungen]
\[
  \Ma_2^2 = \frac{\Ma_1^2 + \frac{2}{\gamma-1}}{\frac{2\gamma}{\gamma-1}\Ma_1^2 - 1}
  \qquad
  \frac{\rho_2}{\rho_1} = \frac{(\gamma+1)\Ma_1^2}{(\gamma-1)\Ma_1^2 + 2}
  \qquad
  \frac{p_2}{p_1} = \frac{2\gamma\Ma_1^2 - (\gamma-1)}{\gamma+1}
\]
Hyperschall-Grenzfall $\Ma_1 \to \infty$:
\[
  \frac{\rho_2}{\rho_1} \to \frac{\gamma+1}{\gamma-1}
    \underset{\gamma=1.4}{=} 6
\]
\end{herleitung}

\begin{figure}[H]
\centering
\begin{tikzpicture}
\begin{axis}[
  width=13cm, height=7.5cm,
  xlabel={Anström-Machzahl $\Ma_1$},
  ylabel={Normiertes Zustandsgrößenverhältnis},
  xmin=1, xmax=30,
  ymin=0, ymax=12,
  grid=both, grid style={draw=gray!25},
  title={Senkrechter Verdichtungsstoß: Zustandsgrößen vs. Machzahl},
  title style={font=\normalsize\bfseries\color{hypblau}},
  legend pos=north west, legend style={font=\small},
]
\addplot[thick, hyprot, domain=1:30, samples=200]
  {(2*1.4*x^2 - (1.4-1))/(1.4+1)/5};
\addlegendentry{$p_2/p_1 \div 5$}
\addplot[thick, hypblau, domain=1:30, samples=200]
  {(1.4+1)*x^2/((1.4-1)*x^2+2)};
\addlegendentry{$\rho_2/\rho_1$}
\addplot[thick, schockgelb!80!black, domain=1:30, samples=200]
  {((2*1.4*x^2-(1.4-1))*((1.4-1)*x^2+2)) / ((1.4+1)^2*x^2) / 3};
\addlegendentry{$T_2/T_1 \div 3$}
\addplot[thick, ablgruen!80!black, domain=1:30, samples=200]
  {sqrt((x^2 + 2/(1.4-1)) / (2*1.4/(1.4-1)*x^2 - 1))};
\addlegendentry{$\Ma_2$}
\addplot[dashed, gray, domain=1:30] {6};
\node[gray, font=\scriptsize] at (axis cs:25,6.3) {Dichtegrenze $=6$};
\addplot[dotted, dunkelgrau] coordinates {(25,0)(25,12)};
\node[dunkelgrau, font=\scriptsize, rotate=90] at (axis cs:24,8) {Shuttle Wiedereintritt};
\end{axis}
\end{tikzpicture}
\caption{Zustandsgrößensprünge am senkrechten Verdichtungsstoß für $\gamma=1.4$.}
\label{fig:normalstoss}
\end{figure}

\subsection{Schiefer Stoß}

\begin{satz}[Stoßwinkelbedingung ($\theta$-$\beta$-$M$)]
\[
  \tan\theta = 2\cot\beta\cdot
    \frac{\Ma_1^2\sin^2\beta - 1}{\Ma_1^2(\gamma + \cos 2\beta) + 2}
\]
\end{satz}

%=======================================================
\section{Aerothermodynamik des Wiedereintritts}
%=======================================================

\subsection{Hochtemperatureffekte in Luft}

\begin{table}[H]
\centering
\caption{Hochtemperatureffekte in Luft bei steigender Temperatur}
\label{tab:hochtemp}
\begin{tabular}{lrl}
\toprule
\textbf{Phänomen} & \textbf{Temp. [K]} & \textbf{Reaktion}\\
\midrule
Vibrationserregung O$_2$        & 800    & Anregung innerer Freiheitsgrade\\
Dissoziation O$_2$              & 2\,000 & $\mathrm{O_2 \to 2\,O}$\\
Dissoziation N$_2$              & 4\,000 & $\mathrm{N_2 \to 2\,N}$\\
Ionisation O, NO                & 9\,000 & $\mathrm{O \to O^+ + e^-}$\\
\bottomrule
\end{tabular}
\end{table}

\subsection{Staupunkttemperatur}

Für den Wiedereintritt aus LEO ($V_\infty = \SI{7800}{m/s}$, $T_\infty = \SI{220}{K}$):
\[
  T_{0,\text{ideal}} = 220 + \frac{(7800)^2}{2\times1005}
    \approx \SI{30\,500}{K}
\]
In der Realität werden durch Dissoziation etwa 65\,\% dieser Energie in chemische
Energie umgewandelt --- ein enormer natürlicher Puffer.

\subsubsection{Wärmefluss am Staupunkt nach Fay-Riddell und Detra-Kemp-Riddell}

\begin{satz}[Fay-Riddell-Formel (vereinfacht)]
\[
  \dot{q}_s = 1{,}83\times10^{-4}\,\frac{1}{\sqrt{R_N}}
    \sqrt{\frac{\rho_\infty}{\rho_{\mathrm{SL}}}}\,
    \left(\frac{V_\infty}{\SI{7924}{m/s}}\right)^{3.15}
    \quad [\si{W/cm^2}]
\]
Schlüsselfolgerung: Größerer Nasenradius $\Rightarrow$ geringerer Wärmefluss ---
Grund für stumpfe Kapselform.
\end{satz}

\begin{figure}[H]
\centering
\begin{tikzpicture}
\begin{axis}[
  width=13cm, height=7cm,
  xlabel={Nasenradius $R_N$ [m]},
  ylabel={Staupunkt-Wärmefluss [W/cm$^2$]},
  xmin=0.01, xmax=2.0,
  ymin=0, ymax=800,
  grid=both, grid style={draw=gray!25},
  title={Wärmefluss am Staupunkt: Detra-Kemp-Riddell ($V=7800$ m/s, $h=60$ km)},
  title style={font=\normalsize\bfseries\color{hypblau}},
  legend pos=north east, legend style={font=\small},
]
\addplot[thick, hyprot, domain=0.05:2.0, samples=100]
  {220/sqrt(x)};
\addlegendentry{$\dot{q}_s \propto R_N^{-1/2}$}
\addplot[only marks, mark=*, mark size=3pt, color=hypblau]
  coordinates {(4.694, 48)};
\addlegendentry{Apollo CM ($R_N=4.69$ m)}
\addplot[only marks, mark=square*, mark size=3pt, color=schockgelb!80!black]
  coordinates {(0.3, 190)};
\addlegendentry{Shuttle Nose ($R_N\approx0.3$ m)}
\draw[dashed, dunkelgrau] (axis cs:0.3,0) -- (axis cs:0.3,700);
\draw[dashed, dunkelgrau] (axis cs:4.694,0) -- (axis cs:4.694,700);
\end{axis}
\end{tikzpicture}
\caption{Wärmefluss am Staupunkt. Die stumpfe Apollo-Kapsel ($R_N = \SI{4.69}{m}$)
reduziert den Wärmefluss drastisch.}
\label{fig:waermefluss}
\end{figure}

\subsection{Grenzschicht und Stoß-Grenzschicht-Interaktion}

Die \emph{adiabate Wandtemperatur} (Recovery Temperature):
\[
  T_{\mathrm{aw}} = T_e \left(1 + r\,\frac{\gamma-1}{2}\Ma_e^2\right)
\]
mit Rückgewinnungsfaktor $r = \Pr^{1/2}$ (laminar) bzw. $\Pr^{1/3}$ (turbulent).

Die \textbf{Stoß-Grenzschicht-Interaktion (SBLI)} erzeugt bei Hyperschall lokale
Wärmefluss-Spitzen bis zum 5-fachen Wandwert, kritisch für Kacheln und Vorderkanten.

%=======================================================
\section{Wärmeschutzsysteme (TPS)}
%=======================================================

\begin{table}[H]
\centering
\caption{Eigenschaften typischer TPS-Materialien}
\label{tab:tps_mat}
\begin{tabular}{lcccll}
\toprule
\textbf{Material} & $T_{\max}$ [°C] & $\rho$ [kg/m$^3$] &
$\lambda$ [W/(mK)] & \textbf{Typ} & \textbf{Anwendung}\\
\midrule
AVCOAT 5026-39    & 3\,300 & 512  & 0.25 & Ablator    & Apollo CM, Orion\\
PICA (Phenolic)   & 3\,000 & 260  & 0.14 & Ablator    & Stardust, Dragon\\
LI-900 (HRSI)     & 1\,260 & 144  & 0.05 & Kachel     & Shuttle Unterseite\\
C/C-SiC           & 2\,000 & 1\,800 & 15  & Struktur. & Hermes (konzept.)\\
\bottomrule
\end{tabular}
\end{table}

Energiebilanz an der Außenwand:
\[
  \rho_s c_s \frac{\partial T}{\partial t}
    = \frac{\partial}{\partial x}\!\left(\lambda_s \frac{\partial T}{\partial x}\right)
    + \dot{q}_{\mathrm{conv}} - \varepsilon\,\sigma\,T_w^4
\]
Strahlungsgleichgewichtstemperatur:
\[
  T_{w,\mathrm{eq}} = \left(\frac{\dot{q}_{\mathrm{conv}}}{\varepsilon\,\sigma}\right)^{1/4}
\]

%=======================================================
\section{Ballistische Flugmechanik beim Wiedereintritt}
%=======================================================

\subsection{Atmosphärisches Eintrittsmodell}

\begin{align}
  m\dot{V} &= -\frac{1}{2}\rho(h) V^2 C_D A - mg\sin\gamma\\
  mV\dot{\gamma} &= \frac{1}{2}\rho(h) V^2 C_L A - mg\cos\gamma
    + \frac{mV^2}{r}\cos\gamma\\
  \dot{h} &= -V\sin\gamma
\end{align}

\begin{herleitung}[Maximalverzögerung beim ballistischen Eintritt]
\[
  \boxed{n_{\max} = \frac{V_e^2\sin|\gamma_e|}{2\,e\,H\,g_0}}
\]
Für Apollo: $V_e = \SI{11}{km/s}$, $\gamma_e = -6.5°$, $H = \SI{7160}{m}$:
$n_{\max} \approx 35.7\,g$
\end{herleitung}

\begin{figure}[H]
\centering
\begin{tikzpicture}
\begin{axis}[
  width=13cm, height=8.5cm,
  xlabel={Fluggeschwindigkeit $V$ [km/s]},
  ylabel={Flughöhe $h$ [km]},
  xmin=0, xmax=8,
  ymin=0, ymax=130,
  grid=both, grid style={draw=gray!25},
  title={Wiedereintritts-Flugbahnen: ballistische vs. gesteuerte Trajektorie},
  title style={font=\normalsize\bfseries\color{hypblau}},
  legend pos=north west, legend style={font=\small},
]
\addplot[thick, hyprot, domain=0:7.8, samples=80]
  {120 - 120*(1 - exp(-((x-7.8)/3.5)^2)) + 10*exp(-((x-1.5)/1.0)^2)};
\addlegendentry{Ballistische Kapsel (Apollo)}
\addplot[thick, hypblau, domain=0:7.8, samples=80]
  {110 - 110/(1+exp(-3*(x-3.5))) + 5*sin(deg(x))};
\addlegendentry{Raumgleiter (Shuttle)}
\addplot[dashed, gray, domain=0:8] {100};
\node[gray, font=\scriptsize] at (axis cs:6,102) {Karman-Linie (100 km)};
\fill[hyprot, opacity=0.06]
  (axis cs:3.5,0) rectangle (axis cs:8,130);
\node[hyprot!80, font=\scriptsize, rotate=90] at (axis cs:7.5,60) {Hyperschall};
\end{axis}
\end{tikzpicture}
\caption{Wiedereintrittstrajektorien. Die ballistische Kapsel folgt einer steilen Bahn;
der Raumgleiter nutzt Auftrieb für einen thermisch schonenderen Eintritt.}
\label{fig:trajektorien}
\end{figure}

%=======================================================
\section{Vergleichende Zusammenfassung aller Antriebskonzepte}
%=======================================================

\begin{figure}[H]
\centering
\begin{tikzpicture}
\begin{loglogaxis}[
  width=13cm, height=10cm,
  xlabel={Spezifischer Impuls $\Isp$ [s]},
  ylabel={Schubkraft $F$ [N]},
  xmin=100, xmax=100000,
  ymin=0.0001, ymax=10000000,
  grid=both,
  grid style={draw=gray!20, line width=0.3pt},
  major grid style={draw=gray!40, line width=0.5pt},
  title={Antriebskonzepte: Schubkraft vs. spezifischer Impuls},
  title style={font=\normalsize\bfseries\color{raumblau}},
  legend style={font=\small, at={(0.02,0.02)}, anchor=south west},
  tick label style={font=\small},
]
\addplot[only marks, mark=square*, mark size=4pt, color=raumblau]
  coordinates {(280, 3e6)};
\addlegendentry{Feststoff (SRB)}
\addplot[only marks, mark=square*, mark size=4pt, color=cyan!70!blue]
  coordinates {(360, 8e5)};
\addlegendentry{LOX/Kerosin (Merlin)}
\addplot[only marks, mark=square*, mark size=4pt, color=triebwerkorange]
  coordinates {(451, 1e6)};
\addlegendentry{LOX/LH$_2$ (SSME/RS-25)}
\addplot[only marks, mark=diamond*, mark size=4pt, color=red!70!black]
  coordinates {(850, 3e5)};
\addlegendentry{NTR (NERVA)}
\addplot[only marks, mark=triangle*, mark size=4pt, color=green!60!black]
  coordinates {(1500, 0.5)};
\addlegendentry{Hall-Effekt (SPT-100)}
\addplot[only marks, mark=triangle*, mark size=4pt, color=purple!70]
  coordinates {(3000, 0.1)};
\addlegendentry{Gitterionen (NSTAR)}
\fill[raumblau, opacity=0.06]
  (axis cs:200,1e5) rectangle (axis cs:500,1e7);
\node[font=\scriptsize, raumblau, opacity=0.8] at (axis cs:320,5e6)
  {Startphase};
\fill[green!70!black, opacity=0.05]
  (axis cs:1000,0.001) rectangle (axis cs:100000,100);
\node[font=\scriptsize, green!60!black, opacity=0.9] at (axis cs:5000,0.01)
  {Tiefraum-/stationär};
\end{loglogaxis}
\end{tikzpicture}
\caption{Vergleich aller Antriebskonzepte: Chemische Triebwerke liefern hohen Schub,
elektrische Antriebe hohen $\Isp$ --- ein fundamentaler Kompromiss.}
\label{fig:vergleich}
\end{figure}

\begin{table}[H]
\centering
\caption{Übersicht aller behandelten Antriebskonzepte}
\label{tab:uebersicht}
\begin{tabular}{lrrll}
\toprule
\textbf{Antrieb} & \textbf{$\Isp$ [s]} & \textbf{$F$ [N]} &
\textbf{Treibmittel} & \textbf{Anwendung}\\
\midrule
Feststoff (SRB)   & 250--280  & $10^5\ldots10^7$ & APCP     & Booster, Start\\
LOX/Kerosin       & 300--360  & $10^4\ldots10^6$ & RP-1/LOX & 1. Stufe\\
LOX/LH$_2$        & 420--460  & $10^4\ldots10^6$ & LH$_2$/LOX& 2./3. Stufe\\
N$_2$O$_4$/UDMH   & 300--320  & $10^3\ldots10^5$ & Hypergol & OMS, RCS\\
NTR (nuklear)     & 800--900  & $10^4\ldots10^5$ & LH$_2$   & Marsreise\\
Hall-Effekt (HET) & 1000--2000 & $0.05\ldots5$   & Xenon    & GEO-Antrieb\\
Gitterionen       & 2000--5000 & $0.001\ldots0.5$& Xenon    & Tiefraumsonde\\
\bottomrule
\end{tabular}
\end{table}

%=======================================================
\section{Schlussbetrachtung}
%=======================================================

Diese Arbeit hat die drei Kerngebiete der modernen Raketenphysik vereint:

\textbf{(1) Antriebslehre:} Die Tsiolkowski-Gleichung $\Dv = v_e\ln(m_0/m_1)$ ist
das fundamentale Designgesetz. Die Lavaldüse wandelt thermische Energie isentrop
in Schub. Chemische Triebwerke ($\Isp=250\ldots460\,\mathrm{s}$) dominieren die
Startphase; elektrische Antriebe ($\Isp$ bis 30\,000\,s) die Tiefraummissionen.

\textbf{(2) Bahnmechanik:} Der Hohmann-Transfer ist das energieoptimale Manöver
zwischen koplanaren Kreisbahnen ($\Dv_{\mathrm{LEO}\to\mathrm{GEO}} = \SI{3.95}{km/s}$).
Die Vis-Viva-Gleichung $v^2 = GM(2/r - 1/a)$ beschreibt alle Kepler-Orbits.
Der \textbf{Wobble-Effekt} --- mechanisch durch exzentrischen Trägheitsring,
Slosh-Kammern und Düsenasymmetrie realisiert --- lässt eine Rakete ab dem
Scheitelpunkt exponentiell zunehmend taumeln, mit lateralen Ablagen von
$100\ldots\SI{400}{m}$, die moderne Abwehrsysteme überfordern.

\textbf{(3) Hyperschallaerodynamik:} Für $\Ma > 12$ versagt das perfekte Gasmodell;
Dissoziation senkt $\gamma_{\mathrm{eff}}$ auf $1{,}1\ldots1{,}2$ und erhöht die
Dichtegrenze auf bis zu 14. Stumpfe Nasengeometrie ($R_N$ groß) ist thermodynamisch
vorteilhaft. SBLI erzeugt lokale Wärmespitzen; TPS-Materialien (Ablator, C/C-SiC)
schützen die Struktur.

\bigskip
\paragraph*{Kernaussagen}
\begin{itemize}[leftmargin=*]
  \item \textbf{Tsiolkowski:}
    $\Dv = v_e\ln(m_0/m_1)$ --- Massenverhältnis begrenzt jede Mission.
  \item \textbf{Hohmann-Transfer:}
    $\Dv_{\mathrm{total}} = (v_{\mathrm{peri}}-v_{c,1}) + (v_{c,2}-v_{\mathrm{apo}})$
    --- energieoptimal, zwei Impulsmanöver.
  \item \textbf{Wobble-Effekt:}
    Wachstumsrate
    $\gamma_\vartheta = \tfrac{m_R R_e^2}{2I_\perp}\cdot\Omega_K$ ---
    passiv/aktiv durch exzentrischen Ring, Slosh-Kammern, Düsenasymmetrie.
    Ablage $\delta r_\perp \sim 100\ldots\SI{400}{m}$.
  \item \textbf{Dichtegrenze:}
    $\rho_2/\rho_1|_{\Ma\to\infty} = (\gamma+1)/(\gamma-1)$ ---
    mit Dissoziation bis Faktor 14.
  \item \textbf{Fay-Riddell:}
    $\dot{q}_s \propto R_N^{-1/2}\,V^3$ --- stumpfe Nase schützt.
  \item \textbf{Mehrstufige Raketen:}
    $\Dv_{\mathrm{ges}} = \sum_k v_{e,k}\ln(m_{0,k}/m_{1,k})$ ---
    sukzessives Abwerfen von Strukturmasse überwindet die Massentyrannei.
\end{itemize}

%=======================================================
\begin{thebibliography}{9}
%=======================================================

\bibitem{tsiolkovsky1903}
K.\,E. Tsiolkovsky:
\textit{Investigation of Outer Space Rocket Devices}.
Nauchnoe Obozrenie (Science Review), 1903.

\bibitem{sutton2017}
G.\,P. Sutton, O. Biblarz:
\textit{Rocket Propulsion Elements}.
9.\,Aufl. Wiley, Hoboken, 2017.

\bibitem{bate1971}
R.\,R. Bate, D.\,D. Mueller, J.\,E. White:
\textit{Fundamentals of Astrodynamics}.
Dover Publications, New York, 1971.

\bibitem{curtis2014}
H.\,D. Curtis:
\textit{Orbital Mechanics for Engineering Students}.
3.\,Aufl. Butterworth-Heinemann, Oxford, 2014.

\bibitem{anderson2006}
J.\,D. Anderson:
\textit{Hypersonic and High Temperature Gas Dynamics}.
2.\,Aufl. AIAA Education Series, Reston, 2006.

\bibitem{fayriddell1958}
J.\,A. Fay, F.\,R. Riddell:
\textit{Theory of Stagnation Point Heat Transfer in Dissociated Air}.
Journal of the Aeronautical Sciences, 25(2):73--85, 1958.

\bibitem{wertz2011}
J.\,R. Wertz, D.\,F. Everett, J.\,J. Puschell (Hrsg.):
\textit{Space Mission Engineering: The New SMAD}.
Microcosm Press, Hawthorne, 2011.

\end{thebibliography}

\end{document}
